\providecommand{\Z}{\mathbb{Z}}
\providecommand{\Q}{\mathbb{Q}}
\providecommand{\C}{\mathbb{C}}
\providecommand{\HH}{\mathbb{H}}
\providecommand{\cO}{\mathcal{O}}
\providecommand{\cG}{\mathcal{G}}
\providecommand{\kBKM}{\kappa_{\mathrm{BKM}}}
\providecommand{\kch}{\kappa_{\mathrm{ch}}}
\providecommand{\kcat}{\kappa_{\mathrm{cat}}}
\providecommand{\kfib}{\kappa_{\mathrm{fiber}}}
\providecommand{\Sp}{\mathrm{Sp}}
\providecommand{\Mp}{\mathrm{Mp}}
\providecommand{\Aut}{\mathrm{Aut}}
\providecommand{\Auts}{\mathrm{Aut}_{s}}
\providecommand{\Fix}{\mathrm{Fix}}
\providecommand{\Kum}{\mathrm{Kum}}
\providecommand{\wt}{\mathrm{wt}}
\providecommand{\BV}{\mathrm{BV}}
\providecommand{\MTW}{\mathrm{MTW}}
\providecommand{\Delfnm}[2]{\Delta_{#1}^{(#2)}}

\section*{Non-CHL Borcherds weights and multiplier rows}
\label{sec:non-chl-borcherds-weights-multiplier-rows}

\paragraph{Two catalogues.}
Two independent sources of Borcherds weights use similar notation.  The
primitive CHL denominator ladder is indexed by the
orders
\[
  N\in \{1,2,3,4,6\}.
\]
For this ladder Gritsenko--Nikulin's construction and Borcherds'
singular-theta weight formula give
\[
  (c_1(0),c_2(0),c_3(0),c_4(0),c_6(0))
  =(10,8,6,4,2),
\]
and therefore
\[
  \kBKM(\Phi_N^{\mathrm{BKM}})
  =\frac{c_N(0)}{2}
  =(5,4,3,2,1).
\]
The genus-two physical denominator convention relevant here records the
square:
\[
  \Phi_N^{\mathrm{phys}}=(\Phi_N^{\mathrm{BKM}})^2,\qquad
  \wt(\Phi_N^{\mathrm{phys}})=2\kBKM(\Phi_N^{\mathrm{BKM}})
  =c_N(0).
\]
Thus the physics weight is the squared-denominator weight, while
\(\kBKM\) is the weight of the primitive BKM denominator.

The Gritsenko--Cl\'ery dd-modular catalogue is different.  Its eight rows
are indexed by triples
\[
\begin{array}{c|c|c|c}
\text{row} & (t,N_{\mathrm{dd}}) & \text{form} & (k,\,c^{\mathrm{dd}}(0))\\
\hline
1 & (1,1) & F_{1,1}=\Delta_5 & (5,10)\\
2 & (2,1) & F_{2,1} & (2,4)\\
3 & (1,2) & F_{1,2} & (3,6)\\
4 & (3,1) & F_{3,1} & (1,2)\\
5 & (1,3) & F_{1,3} & (2,4)\\
6 & (4,1) & F_{4,1} & (\tfrac12,1)\\
7 & (1,4) & F_{1,4} & (\tfrac32,3)\\
8 & (2,2) & F_{2,2} & (1,2)
\end{array}
\]
with \(c^{\mathrm{dd}}(0)=2k\).  These rows live on their own
paramodular groups \(\Gamma_t(N_{\mathrm{dd}})\), with their own
characters and divisor data.  Their constants form the tuple
\[
  c^{\mathrm{dd}}(0)=(10,4,6,2,4,1,3,2),
\]
with half-integral weights \(1/2\) and \(3/2\) appearing in the rows
\((4,1)\) and \((1,4)\).  The common row
\(F_{1,1}=\Delta_5\) is the sole automatic overlap with the primitive
CHL ladder; every further overlap requires a theorem identifying the
row and the normalisation.

\paragraph{\texorpdfstring{\(\Delta_5\)}{Delta5} and
\texorpdfstring{\(\Phi_{10}\)}{Phi10}.}
The primitive \(N=1\) BKM denominator uses the Eichler--Zagier weak
Jacobi form \(\phi_{0,1}\), with \(q^0\)-head
\[
  \phi_{0,1}|_{q^0}=y+10+y^{-1}.
\]
The K3 elliptic genus is \(Z_{K3}=2\phi_{0,1}\).  Using
\(2\phi_{0,1}\) doubles every Borcherds exponent and lands on the
Igusa-square convention
\[
  \Phi_{10}^{\mathrm{un}}=\Delta_5^2 .
\]
The primitive denominator has weight \(5\); the square has weight
\(10\).

\paragraph{\texorpdfstring{\(K3\times E\)}{K3xE} invariants.}
On the compact product row the invariants are
\[
  \kcat(K3\times E)=0,\qquad
  \kch(K3\times E)=0,\qquad
  \kch^{\mathrm{Heis}}(K3\times E)=3,
\]
and
\[
  \kBKM(\Delta_5)=5,\qquad
  \kfib(K3\times E)=24.
\]
They come from distinct constructions: compact Hodge supertrace,
Heisenberg--Mukai specialisation, Borcherds weight, and Mukai-lattice
rank of the K3 fibre.

\paragraph{Hall--Drinfeld and Yangian branches.}
The \(K3\times E\) Borcherds object is the Hall--Drinfeld double or
compact Hall branch attached to the positive Hall half.  The Mukai
self-mirror K3 Yangian branch is a separate stable-envelope and current
algebra construction.  The BKM denominator identity for the
Hall--Drinfeld branch comes from the Borcherds product and its Weyl
denominator data.

\paragraph{Borcherds weight theorem.}
Let \(\psi\) be the Jacobi input for a Borcherds product on an even
lattice of signature \((2,n)\), with constant coefficient \(c(0)\).
Borcherds' weight formula gives
\[
  \wt(\mathrm{Bor}(\psi))=\frac{c(0)}{2}.
\]
This is a theorem about the automorphic product.  A BKM denominator
identity requires additional root-lattice, Weyl-vector, and
simple-imaginary-root data.  On the primitive CHL ladder these data are
provided by Gritsenko--Nikulin.  On a multiplier row of the
Gritsenko--Cl\'ery catalogue they must be supplied row by row.

\paragraph{Mukai orders outside the CHL denominator ladder.}
Symplectic automorphisms of K3 surfaces exist for
\[
  N\in \{1,2,3,4,5,6,7,8\},
\]
while elliptic curves have automorphism orders only \(1,2,3,4,6\).
For \(N=5,7,8\), the ordinary free quotient
\((K3\times E)/\Z_N\) lacks the elliptic automorphism factor required
by the primitive CHL denominator ladder.
The order-indexed Borcherds lane attached to the K3-twined Jacobi input
has
\[
  (c_5(0),c_7(0),c_8(0))=(4,2,2),\qquad
  \kBKM=(2,1,1).
\]
These three numbers belong to the order-indexed Nikulin lane.  The
half-integral multiplier rows \((4,1;\tfrac12)\) and
\((1,4;\tfrac32)\) belong to the dd-catalogue.  Numerical equality of
weights gives no automorphic-product isomorphism.

\subsection*{The order \(N=5\)}

\paragraph{Automorphic datum.}
For the K3-twined order-indexed Borcherds product the constant is
\[
  c_5(0)=4,\qquad \kBKM(\Phi_5^{\mathrm{ord}})=2.
\]
This is the Borcherds-weight reading.  It is compatible with the
Jatkar--Sen value \(24/(5+1)-2=2\).  It belongs to neither the
primitive CHL denominator ladder nor the dd-row \(F_{4,1}\) of weight
\(\tfrac12\).

\paragraph{Geometric host.}
An order-five symplectic K3 automorphism exists, with four fixed points
by Nikulin and Mukai.  Elliptic curves have no order-five automorphism
fixing the origin, so a \(K3\times E\) CHL quotient of order five is
absent.  A Borcea--Voisin construction supplies a replacement
order-five host after choosing a K3 surface \(S_5\), an antisymplectic
involution \(\iota_S\) commuting with the order-five symplectic
automorphism, and the elliptic involution \(\iota_E(z)=-z\):
\[
  X_5^{\BV}:=(S_5\times E)/(\iota_S\times\iota_E).
\]
After crepant resolution this gives a Calabi--Yau threefold
\(\widehat X_5^{\BV}\) in the usual Borcea--Voisin sense.  Its
categorical Hodge supertrace is
\[
  \kcat(\widehat X_5^{\BV})
  =\chi(\cO_{\widehat X_5^{\BV}})
  =h^{0,0}-h^{0,1}+h^{0,2}-h^{0,3}=1-0+0-1=0.
\]
This Hodge calculation leaves the chiral invariant \(\kch\)
undetermined.  It requires an explicit \(\Phi\)-image of
\(D^b\mathrm{Coh}(\widehat X_5^{\BV})\), or an equivalent
factorisation-algebra model, together with the chosen order-five
equivariance.

\paragraph{Denominator claim.}
The proposed Borcea--Voisin DT denominator remains an unproved
identification:
\[
  Z^{\widehat X_5^{\BV}}_{\mathrm{red},L}
  \stackrel{\mathrm{conj}}{=}
  -\,C_5\bigl(\Phi_5^{\mathrm{ord}}\bigr)^{-2},
\]
for a correction \(C_5\) fixed by the chamber and curve class.  A fourth
root or a metaplectic square root requires naming the
multiplier row and proving that the DT denominator is a power of that
row.  The CHL ladder alone supplies no such identification.

\subsection*{The order \(N=7\)}

\paragraph{Automorphic datum.}
For the order-seven K3-twined Borcherds product,
\[
  c_7(0)=2,\qquad \kBKM(\Phi_7^{\mathrm{ord}})=1.
\]
This is separate from the dd-catalogue rows.  An identification of an
order-seven K3 twining with a quarter-weight paramodular form would
require a new overlap theorem.

\paragraph{Geometric obstruction.}
An order-seven symplectic K3 automorphism has three fixed points.
Elliptic curves have no order-seven automorphism, so the ordinary
\((K3\times E)/\Z_7\) CHL construction is absent.  A quotient of
\(S_7\times E\) by an order-seven action has fixed local models in
place of a free Calabi--Yau threefold quotient.  Any singular or
gerbe-twisted host must specify:
\[
  X_7,\qquad [\cG_7]\in H^3(X_7,\Z)_{\mathrm{tors}},\qquad
  \Phi(D^b\mathrm{Coh}(X_7,\cG_7)).
\]
Only after these data are fixed do the invariants
\(\kcat(X_7,\cG_7)\) and \(\kch(X_7,\cG_7)\) have a defined meaning.
In particular, a rational expression such as \(6/7\) cannot be the
holomorphic Euler characteristic of a smooth Calabi--Yau threefold.

\paragraph{Proof obligation.}
The order-seven Borcherds-weight identity is automorphic.  A BKM or BPS
interpretation requires a twisted virtual class, a Cheeger--Simons
character matching the automorphic multiplier, and coefficient matching
between the resulting partition function and the Borcherds product.

\subsection*{The order \(N=8\)}

\paragraph{Automorphic datum.}
For the order-eight K3-twined Borcherds product,
\[
  c_8(0)=2,\qquad \kBKM(\Phi_8^{\mathrm{ord}})=1.
\]
Zero-weight language belongs only to a separately named
automorphic-product degeneration; the order-eight K3-twined Borcherds
constant is \(2\).

\paragraph{Geometric host.}
Mongardi--Tari--Wandel construct order-eight symplectic automorphisms on
generalised Kummer varieties.  This gives a hyperk\"ahler host outside
the Calabi--Yau-threefold category.  Therefore it cannot be inserted
directly into the \(d=3\) CY-to-chiral functor without a
dimensional-reduction theorem.
For a chosen quotient or stack \(X_8^{\MTW}\), the invariants
\[
  \kcat(X_8^{\MTW})=\chi(\cO_{X_8^{\MTW}}),\qquad
  \kch(\Phi(D^b\mathrm{Coh}(X_8^{\MTW})))
\]
must be computed from that host.  The \(K3\times E\) values and the
number \(1\) in the order-indexed Borcherds lane are comparison data
only.

\paragraph{Boundary interpretation.}
A WZW or affine Kac--Moody degeneration is a physical heuristic unless
one supplies a vertex-algebra character calculation and identifies it
with the named automorphic product.  The automorphic theorem supplies
\(\kBKM=1\) for the order-indexed product; the algebraic
degeneration is a separate construction.

\subsection*{Compact table}

\[
\renewcommand{\arraystretch}{1.15}
\begin{array}{c|c|c|c|c}
\text{source} & \text{index} & c(0) & \kBKM=c(0)/2 & \text{geometric scope}\\
\hline
\text{primitive CHL ladder} & 1,2,3,4,6 & 10,8,6,4,2
  & 5,4,3,2,1 & (K3\times E)/\Z_N\\
\text{order-indexed non-CHL lane} & 5,7,8 & 4,2,2
  & 2,1,1 & \text{new host required}\\
\text{dd-catalogue row} & (4,1) & 1
  & \tfrac12 & \text{multiplier row}\\
\text{dd-catalogue row} & (1,4) & 3
  & \tfrac32 & \text{multiplier row}
\end{array}
\]

\subsection*{Proof obligations}

\begin{enumerate}
\item Construct a reduced DT or stable-pairs theory for the proposed
order-five Borcea--Voisin host and prove its modular denominator.
\item Define the order-seven gerbe-twisted host and compute
\(\kcat\) and \(\kch\) from that host, with \(K3\times E\) serving only
as comparison datum.
\item Prove any dimensional reduction from an order-eight generalised
Kummer host to a CY$_3$ chiral algebra before assigning a \(d=3\)
\(\kch\).
\item For every multiplier row, name the row \((t,N_{\mathrm{dd}};k)\),
its constant \(c^{\mathrm{dd}}(0)\), its cover or character, and the
root datum of the proposed denominator algebra.
\end{enumerate}

\paragraph{Sources.}
Gritsenko--Nikulin, \emph{Automorphic forms and Lorentzian
Kac--Moody algebras II}, 1995, Part II, Theorem 2.1; Borcherds,
\emph{Automorphic forms with singularities on Grassmannians}, 1998,
Theorem 13.3; Gritsenko--Cl\'ery, \emph{The Siegel modular forms of
genus two with the simplest divisor}, 2008/2013, Theorems 1.1 and 3.1;
Nikulin, \emph{Finite groups of automorphisms of K\"ahlerian K3
surfaces}, 1980, Theorem 4.5; Mukai, \emph{Finite groups of
automorphisms of K3 surfaces}, 1988; Borcea 1997; Voisin 1993;
Mongardi--Tari--Wandel 2018.
